\appendix
\chapter{Reproduction Guide}
\label{ch:appendix}

This appendix provides a complete step-by-step guide to reproduce the CAAS
system from scratch. All commands are run from the repository root unless
noted otherwise.

\section{Prerequisites}

\begin{itemize}
  \item Python 3.11 (tested; 3.10+ should work)
  \item Git
  \item Docker (for containerised deployment)
  \item Terraform $\geq$ 1.6 (for AWS provisioning)
  \item AWS account with programmatic access (S3 + EC2; free tier sufficient for testing)
  \item NASA FIRMS API key --- register free at \url{https://firms.modaps.eosdis.nasa.gov/api/}
  \item GitHub account with Actions enabled
\end{itemize}

\section{Step 1 — Clone and Environment Setup}

\begin{lstlisting}[language=bash]
git clone https://github.com/supanut-k/caas.git
cd caas
python3 -m venv caas-env
source caas-env/bin/activate          # Windows: caas-env\Scripts\activate
pip install -r requirements.txt
\end{lstlisting}

\section{Step 2 — Configure Environment Variables}

\begin{lstlisting}[language=bash]
cp .env.example .env
\end{lstlisting}

Edit \texttt{.env} and fill in:

\begin{lstlisting}
FIRMS_MAP_KEY=<your_firms_api_key>
AWS_ACCESS_KEY_ID=<your_aws_key>
AWS_SECRET_ACCESS_KEY=<your_aws_secret>
AWS_REGION=ap-southeast-1
S3_BUCKET_NAME=caas-mlops-bucket
MLFLOW_TRACKING_URI=http://localhost:5000
GITHUB_TOKEN=<your_github_pat>
GITHUB_REPOSITORY=supanut-k/caas
\end{lstlisting}

\section{Step 3 — Provision AWS Infrastructure with Terraform}

\begin{lstlisting}[language=bash]
cd infra/
terraform init
terraform plan                        # Review what will be created
terraform apply                       # Creates S3 bucket, IAM role, EC2 t3.small
\end{lstlisting}

After apply, note the outputs:
\begin{lstlisting}
api_url     = "http://<ec2-public-ip>:8000"
mlflow_url  = "http://<ec2-public-ip>:5000"
\end{lstlisting}

Update \texttt{MLFLOW\_TRACKING\_URI} in \texttt{.env} to the EC2 MLflow URL.

\section{Step 4 — Historical Data Backfill}

Fetch all historical data from 2011 to the present day:

\begin{lstlisting}[language=bash]
# Fetch PM2.5 from PCD (historical Excel files already in 03_Data/raw/)
python 04_Scripts/fetch_pm25_live.py

# Fetch ERA5 weather via Open-Meteo API
python 04_Scripts/fetch_weather.py

# Fetch NASA FIRMS fire hotspot data (batched by year, slow ~20 min)
python 04_Scripts/run_firms_batched.py

# Build the 45-feature matrix
python 04_Scripts/build_features.py
\end{lstlisting}

The resulting \texttt{03\_Data/processed/features.csv} should contain
$\geq$ 5,000 daily rows spanning 2011--present.

\section{Step 5 — Train Models and Register with MLflow}

Start MLflow tracking server first:

\begin{lstlisting}[language=bash]
mlflow server --host 0.0.0.0 --port 5000 &
\end{lstlisting}

Train both model families:

\begin{lstlisting}[language=bash]
# XGBoost (t+1, t+3, t+7) — ~2 min on modern CPU
python 04_Scripts/train_xgboost.py

# LSTM (t+1, t+3, t+7) — ~60 min on CPU, ~15 min on GPU
python 04_Scripts/train_lstm.py
\end{lstlisting}

Inspect runs at \url{http://localhost:5000}. XGBoost models are automatically
promoted to \texttt{Production} stage by \texttt{train\_xgboost.py}.
LSTM models are logged as \texttt{Staging} for comparison only.

\section{Step 6 — Run Analysis Scripts (Optional)}

\begin{lstlisting}[language=bash]
# Scenario C: alert threshold optimisation + PR curves
python 04_Scripts/scenario_c_threshold.py

# SHAP feature importance across all horizons
python 04_Scripts/shap_analysis.py
\end{lstlisting}

Figures are saved to \texttt{03\_Data/results/} and
\texttt{07\_Final/report/figures/}.

\section{Step 7 — Run Inference and Upload to S3}

\begin{lstlisting}[language=bash]
# Generate latest forecasts from the feature matrix
python 04_Scripts/run_inference.py

# Upload data, models, and results to S3
python 04_Scripts/upload_to_s3.py
# Or for a dry run to see what would be uploaded:
python 04_Scripts/upload_to_s3.py --dry-run
\end{lstlisting}

\section{Step 8 — Build and Deploy Docker Container}

\begin{lstlisting}[language=bash]
# Build the slim serving image (~300 MB)
docker build -t caas-api:latest .

# Test locally
docker run -p 8000:8000 \
  -e AWS_REGION=ap-southeast-1 \
  -e S3_BUCKET_NAME=caas-mlops-bucket \
  caas-api:latest

# Verify API is running
curl http://localhost:8000/health
curl http://localhost:8000/forecast
\end{lstlisting}

For production deployment, SSH into the EC2 instance (provisioned in Step 3)
and run the same Docker commands. The Terraform user-data script handles this
automatically on first launch.

\section{Step 9 — Enable CI/CD Automation}

Add the following GitHub Actions secrets to the repository
(\texttt{Settings -> Secrets and variables -> Actions}):

\begin{itemize}
  \item \texttt{FIRMS\_MAP\_KEY} --- NASA FIRMS API key
  \item \texttt{AWS\_ACCESS\_KEY\_ID} and \texttt{AWS\_SECRET\_ACCESS\_KEY}
  \item \texttt{S3\_BUCKET\_NAME} --- \texttt{caas-mlops-bucket}
  \item \texttt{EC2\_HOST} --- public IP of the EC2 instance
  \item \texttt{EC2\_SSH\_KEY} --- private SSH key content
  \item \texttt{MLFLOW\_TRACKING\_URI} --- \texttt{http://<ec2-ip>:5000}
  \item \texttt{GITHUB\_TOKEN} --- automatically available in Actions
\end{itemize}

Once secrets are set, the three workflows activate automatically:

\begin{itemize}
  \item \textbf{test.yml} --- runs on every push to \texttt{main}
  \item \textbf{daily\_pipeline.yml} --- runs at 06:00 Bangkok time (23:00 UTC)
  \item \textbf{retrain.yml} --- triggered by drift alert or December 1 schedule
\end{itemize}

\section{Step 10 — Launch the Streamlit Dashboard}

\begin{lstlisting}[language=bash]
cd 04_Scripts/serve
streamlit run dashboard.py
# Opens at http://localhost:8501
\end{lstlisting}

The dashboard reads the latest forecast from
\texttt{03\_Data/results/latest\_forecast.json} (updated daily by the pipeline).
If the file is absent, it falls back to a demo mode with synthetic data.

\section{Step 11 — Run the Test Suite}

\begin{lstlisting}[language=bash]
pytest                    # Runs all 38 tests with -v --tb=short
pytest tests/test_api.py  # API endpoint tests only
pytest tests/test_models.py  # Model file + inference tests
\end{lstlisting}

All 38 tests should pass. Tests that require model files will skip gracefully
if models are not yet trained.

% ------------------------------------------------------------------

\chapter{Full Metric Tables}
\label{app:metrics}

\section{XGBoost — Validation and Test Metrics}

\begin{table}[H]
\centering
\caption{XGBoost Full Metrics (Validation and Test Sets)}
\begin{tabular}{llrrrr}
\toprule
\textbf{Split} & \textbf{Horizon} & \textbf{MAE} & \textbf{RMSE} & \textbf{R²} & \textbf{AUROC} \\
\midrule
Validation & t+1 & 5.74  & 9.21  & 0.799 & 0.988 \\
Validation & t+3 & 7.85  & 12.96 & 0.701 & 0.963 \\
Validation & t+7 & 9.50  & 15.92 & 0.579 & 0.946 \\
\midrule
Test       & t+1 & 5.31  & 9.48  & 0.824 & 0.991 \\
Test       & t+3 & 6.91  & 11.62 & 0.736 & 0.976 \\
Test       & t+7 & 8.34  & 14.13 & 0.612 & 0.959 \\
\bottomrule
\end{tabular}
\end{table}

\section{LSTM — Validation and Test Metrics}

\begin{table}[H]
\centering
\caption{LSTM Full Metrics (Validation and Test Sets)}
\begin{tabular}{llrrrr}
\toprule
\textbf{Split} & \textbf{Horizon} & \textbf{MAE} & \textbf{RMSE} & \textbf{R²} & \textbf{AUROC} \\
\midrule
Validation & t+1 & 7.12  & 11.45 & 0.730 & 0.978 \\
Validation & t+3 & 9.45  & 16.20 & 0.512 & 0.914 \\
Validation & t+7 & 9.87  & 16.78 & 0.543 & 0.932 \\
\midrule
Test       & t+1 & 6.67  & 11.64 & 0.753 & 0.984 \\
Test       & t+3 & 8.87  & 15.89 & 0.540 & 0.927 \\
Test       & t+7 & 8.06  & 14.49 & 0.608 & 0.961 \\
\bottomrule
\end{tabular}
\end{table}

\section{Alert Classification Metrics (XGBoost, 50 µg/m³ threshold)}

\begin{table}[H]
\centering
\caption{Binary Alert Classification Performance on Test Set}
\begin{tabular}{lrrrr}
\toprule
\textbf{Horizon} & \textbf{Precision} & \textbf{Recall} & \textbf{F1} & \textbf{AUROC} \\
\midrule
t+1 & 0.769 & 0.924 & 0.839 & 0.991 \\
t+3 & 0.762 & 0.815 & 0.787 & 0.976 \\
t+7 & 0.667 & 0.702 & 0.684 & 0.959 \\
\bottomrule
\end{tabular}
\end{table}

\section{SHAP Mean Absolute Values — Top 5 Features per Horizon}

\begin{table}[H]
\centering
\caption{Top-5 SHAP Features by Horizon (Mean $|$SHAP$|$ on Test Set)}
\begin{tabular}{clr}
\toprule
\textbf{Horizon} & \textbf{Feature} & \textbf{Mean $|$SHAP$|$} \\
\midrule
t+1 & \texttt{pm25\_lag1}        & 6.85 \\
t+1 & \texttt{alert\_lag1}       & 2.14 \\
t+1 & \texttt{pm25\_roll7\_mean} & 1.93 \\
t+1 & \texttt{hotspot\_50km}     & 1.21 \\
t+1 & \texttt{pm25\_lag3}        & 0.98 \\
\midrule
t+3 & \texttt{hotspot\_14d\_roll} & 4.63 \\
t+3 & \texttt{pm25\_lag3}         & 3.87 \\
t+3 & \texttt{pm25\_roll7\_mean}  & 2.95 \\
t+3 & \texttt{sin\_month}         & 1.74 \\
t+3 & \texttt{hotspot\_7d\_roll}  & 1.52 \\
\midrule
t+7 & \texttt{sin\_month}         & 4.63 \\
t+7 & \texttt{week}               & 2.15 \\
t+7 & \texttt{hotspot\_14d\_roll} & 2.01 \\
t+7 & \texttt{pm25\_roll14\_mean} & 1.88 \\
t+7 & \texttt{cos\_month}         & 1.34 \\
\bottomrule
\end{tabular}
\end{table}
